\documentclass[a4paper,12pt]{article}


\usepackage{parskip}
\usepackage[pdftex]{graphicx}

\usepackage[utf8]{inputenc}
\usepackage[T1]{fontenc}
\usepackage[british]{babel}

\usepackage{amsmath}
\usepackage{amsthm}

\usepackage{typearea}

\usepackage[toc,page]{appendix}
\usepackage{float}
% Acronyms
\usepackage{arydshln}





% Divide your \newtheorem commands into groups
% Comment out the necessary commands

\theoremstyle{plain}
\newtheorem{theorem}{Theorem}
%\newtheorem{lemma}{Lemma}
\newtheorem{corollary}{Corollary}
\newtheorem{proposition}{Proposition}

\theoremstyle{definition}
\newtheorem{definition}{Definition}
%\newtheorem{condition}{Condition}
%\newtheorem{problem}{Problem}
%\newtheorem{example}{Example}

\theoremstyle{remark}
%\newtheorem{remark}{Remark}
%\newtheorem{note}{Note}
%\newtheorem{claim}{Claim}

\usepackage{newtxtext}
\usepackage{newtxmath}

\begin{document}


\begin{center}
\huge Algebraic Geometry\\
\vspace{0.2cm}
\Large Solutions to selected problems for assignment 3 \\  \vspace{0.5cm}
Mara Kali\v{c}anin Dimitrov
\vspace{0.5cm}
\end{center}



\subsection*{Problem 8.1}

Let
\[
P=[a_0:\cdots:a_n]\in P^n.
\]
We need to find a homogeneous ideal $I\subseteq k[x_0,\dots,x_n]$ such that
\[
V(I)=\{P\}.
\]

\medskip

Let's recall the relevant definitions from Olof's book:
\begin{itemize}
\item \textbf{Definition 8.1}. Projective $n$-space $P^n$ is the set of equivalence classes
\[
[a_0:\cdots:a_n]
\]
with $(a_0,\dots,a_n)\neq (0,\dots,0)$, where two nonzero tuples determine the same point if and only if they differ by multiplication by a nonzero scalar.

\item \textbf{Definition 8.3}. A homogeneous ideal is an ideal generated by homogeneous elements.

\item \textbf{Definition 8.4}. A point $Q\in P^n$ is a zero of a homogeneous polynomial $f$ if
\[
f(b_0,\dots,b_n)=0
\]
for any choice of homogeneous coordinates
\[
[b_0:\cdots:b_n]
\]
for $Q$.

\item \textbf{Definition 8.5}. If $S$ is a set of homogeneous polynomials, then
\[
V(S)=\{Q\in P^n \mid f(Q)=0 \text{ for all } f\in S\}.
\]
\end{itemize}

\medskip

Since $P$ is a point of $P^n$, not all $a_0,\dots,a_n$ are zero. Choose an index $s$ such that
\[
a_s\neq 0.
\]
Now define
\[
I=\bigl(a_jx_i-a_ix_j \mid 0\le i,j\le n\bigr)\subseteq k[x_0,\dots,x_n].
\] 

Each generator
\[
a_jx_i-a_ix_j
\]
is homogeneous of degree $1$. Therefore, by Definition 8.3, $I$ is a homogeneous ideal.

\medskip

Next we show that $V(I)=\{P\}$.

First, $P\in V(I)$. Indeed,
\[
[a_0:\cdots:a_n]
\]
is a choice of homogeneous coordinates for $P$, and for every pair $i,j$ we have
\[
(a_jx_i-a_ix_j)(a_0,\dots,a_n)=a_ja_i-a_ia_j=0.
\]
Thus every generator of $I$ vanishes on these homogeneous coordinates of $P$, so by Definitions 8.4 and 8.5 we have
\[
P\in V(I).
\]

\medskip

Next, let
$
Q=[b_0:\cdots:b_n]\in V(I).$
Since
$
a_sx_i-a_ix_s\in I
$
for every $i$, Definition 8.5 gives
\[
a_sb_i-a_ib_s=0
\qquad\text{for all } i=0,\dots,n.
\]
These equations express that the coordinate vector
$
(b_0,\dots,b_n)
$
is proportional to
$
(a_0,\dots,a_n),
$
and we now make this explicit.

Next, show that $b_s\neq 0$. If $b_s=0$, then
\[
a_sb_i=0
\qquad\text{for all } i.
\]
Since $a_s\neq 0$, it follows that $b_i=0$ for all $i$, which is impossible because $Q$ is a point of $P^n$.
Hence
\[
b_s\neq 0.
\]

Now set
\[
\lambda=\frac{b_s}{a_s}\in k^\times.
\]
Then for every $i$,
\[
a_sb_i=a_ib_s=a_i(\lambda a_s).
\]
Since $a_s\neq 0$, we may divide by $a_s$ and obtain
\[
b_i=\lambda a_i
\qquad\text{for all } i=0,\dots,n.
\]
Therefore
\[
(b_0,\dots,b_n)=\lambda(a_0,\dots,a_n).
\]
By Definition 8.1, this means that
\[
Q=[b_0:\cdots:b_n]=[a_0:\cdots:a_n]=P.
\]

Thus every point of $V(I)$ is equal to $P$. Since we already showed that $P\in V(I)$, we can conclude that
\[
V(I)=\{P\}.
\]

Therefore, we showed that 
\[
I=\bigl(a_jx_i-a_ix_j \mid 0\le i,j\le n\bigr)
\] 
is the homogeneous ideal such that \[
V(I)=\{P\}.
\]



\subsection*{Problem 8.5}

Let $V_1,\dots,V_n$ be hyperplanes in $P^n$. We need to show that
\[
\bigcap_{i=1}^n V_i\neq \emptyset.
\]

\medskip

To show this, we want to reduce this to a homogeneous linear system with $n$ equations in $n+1$ variables.

For each $i=1,\dots,n$, let
\[
V_i=V(\ell_i)
\]
where $\ell_i$ is a nonzero homogeneous polynomial of degree $1$. Write
\[
\ell_i(x_0,\dots,x_n)=a_{i0}x_0+\cdots+a_{in}x_n.
\]
Thus we need to show that there exists a point
\[
P=[b_0:\cdots:b_n]\in P^n
\]
such that
\[
\ell_i(b_0,\dots,b_n)=0
\qquad\text{for all } i=1,\dots,n.
\]

Equivalently, we need to show that the homogeneous linear system
\[
\sum_{j=0}^n a_{ij}x_j=0,
\qquad i=1,\dots,n,
\]
has a nonzero solution in $k^{n+1}$. To prove this, we will use induction on $n$.
 
 
 For $n=1$, we have one nonzero linear equation in two variables,
\[
a_{10}x_0+a_{11}x_1=0,
\]
where $(a_{10},a_{11})\neq (0,0)$. Consider the vector
\[
(-a_{11},a_{10}).
\]
This vector is nonzero, and
\[
a_{10}(-a_{11})+a_{11}a_{10}=0.
\]
Hence, the system has a nonzero solution when $n=1$.

\medskip
Next, assume that any homogeneous system of $n-1$ linear equations in $n$ variables has a nonzero solution. Now consider a homogeneous system of $n$ linear equations in $n+1$ variables:
\[
\sum_{j=0}^n a_{ij}x_j=0,
\qquad i=1,\dots,n.
\]

We consider two cases.

1. Let $a_{in}=0$ for all $i=1,\dots,n$.
Then the variable $x_n$ does not appear in any equation. Therefore
\[
(0,\dots,0,1)
\]
is a nonzero solution of the whole system.

2. $a_{in}\neq 0$ for at least one $i$.
After renumbering the equations if necessary, we may assume that
\[
a_{1n}\neq 0.
\]
Then the first equation gives
\[
a_{10}x_0+\cdots+a_{1,n-1}x_{n-1}+a_{1n}x_n=0,
\]
and since $a_{1n}\neq 0$ and $k$ is a field, we may divide by $a_{1n}$ to obtain
\[
x_n=-a_{1n}^{-1}\sum_{j=0}^{n-1}a_{1j}x_j.
\]
Substituting this into the remaining equations, we obtain
\[
\sum_{j=0}^{n-1}\left(a_{ij}-a_{in}a_{1n}^{-1}a_{1j}\right)x_j=0,
\qquad i=2,\dots,n.
\]
Thus we get a homogeneous system of $n-1$ linear equations in the $n$ variables
\[
x_0,\dots,x_{n-1}.
\]
By the induction hypothesis, this reduced system has a nonzero solution
\[
(b_0,\dots,b_{n-1})\in k^n.
\]
Now define
\[
b_n=-a_{1n}^{-1}\sum_{j=0}^{n-1}a_{1j}b_j.
\]
Then $(b_0,\dots,b_n)$ satisfies the first equation by construction, and it satisfies the remaining equations because $(b_0,\dots,b_{n-1})$ satisfies the reduced system. Since $(b_0,\dots,b_{n-1})$ is nonzero, the vector
\[
(b_0,\dots,b_n)
\]
is also nonzero.

Therefore, the original homogeneous system has a nonzero solution.

\medskip

By induction, the system
\[
\sum_{j=0}^n a_{ij}x_j=0,
\qquad i=1,\dots,n,
\]
always has a nonzero solution
\[
(b_0,\dots,b_n)\in k^{n+1}.
\]
By Definition 8.1, this determines a point
\[
P=[b_0:\cdots:b_n]\in P^n.
\]
Since
\[
\ell_i(b_0,\dots,b_n)=0
\qquad\text{for all } i,
\]
Definition 8.4 shows that $P$ is a zero of each $\ell_i$. Hence by Definition 8.5,
\[
P\in V(\ell_i)=V_i
\qquad\text{for all } i=1,\dots,n.
\]
Therefore
\[
P\in \bigcap_{i=1}^n V_i,
\]
and so
\[
\bigcap_{i=1}^n V_i\neq \emptyset.
\]



\subsection*{Problem 9.1}

Let
\[
N=\frac{d(d+3)}{2}
\]
and let $P_1,\dots,P_N$ be points in $P^2$. We need to show that there is at least one curve of degree $d$ passing through all of the points $P_i$.

\medskip

Let's recall the relevant definitions and results from Olof's book:
\begin{itemize}
\item \textbf{Definition 9.1}. A projective plane curve is an equivalence class of nonzero homogeneous polynomials in $k[x_0,x_1,x_2]$ under multiplication by a nonzero scalar.

\item At the beginning of \S 9.2, after choosing a monomial basis
\[
M_0,\dots,M_N
\]
of the homogeneous polynomials of degree $d$, the set $C_d$ of curves of degree $d$ is identified with the projective space
\[
P^N,
\qquad
N=\frac{d(d+3)}{2}.
\]

\item \textbf{Lemma 9.6}. For any point $P\in P^2$, the set of curves of degree $d$ containing $P$ forms a hyperplane.

\item By Problem 8.5, if $V_1,\dots,V_n$ are hyperplanes in $P^n$, then
\[
\bigcap_{i=1}^n V_i\neq \emptyset.
\]
\end{itemize}

\medskip

Let
\[
A=k[x_0,x_1,x_2],
\]
and let $A_d$ be the vector space of homogeneous polynomials of degree $d$.
Choose the monomials of degree $d$
\[
M_0,\dots,M_N
\]
as a basis of $A_d$, where
\[
N=\frac{(d+1)(d+2)}{2}-1=\frac{d(d+3)}{2}.
\]
Then every nonzero $f\in A_d$ can be written uniquely as
\[
f=a_0M_0+\cdots+a_NM_N.
\]

By Definition 9.1, two nonzero homogeneous polynomials define the same curve if and only if they differ by multiplication by a nonzero scalar. Therefore the parameter space $C_d$ of projective plane curves of degree $d$ can be identified with the projective space
$P^N.$

For each $i=1,\dots,N$, let
\[
V_i=\{\, C\in C_d \mid P_i\in C \,\}.
\]
Thus $V_i$ is the set of degree $d$ curves passing through the point $P_i$.

By Lemma 9.6, each condition
\[
P_i\in C
\]
defines a hyperplane in this same projective space
$P^N.$
Therefore Problem 8.5 applies with

$n=N.$
Since we have exactly $N$ such hyperplanes, we obtain
\[
\bigcap_{i=1}^N V_i\neq \emptyset.
\]
Next, choose a curve
$C\in \bigcap_{i=1}^N V_i.$
Then
\[
P_i\in C
\qquad\text{for all } i=1,\dots,N.
\]
Therefore there is at least one curve of degree $d$ passing through all of the points
\[
P_1,\dots,P_N.
\]



\subsection*{Problem 9.2}

Let $P_1,\dots,P_r$ be points in $P^2$ in general position.

We need to do the following
\begin{itemize}
\item[a)] show that there is a unique line passing through $P_1$ and $P_2$,
\item[b)] show that there is a unique conic passing through $P_1,\dots,P_5$,
\item[c)] show that there is a unique cubic passing through $P_1,\dots,P_9$,
\item[d)] determine how many points in general position are required to specify a quartic.
\end{itemize}

\medskip

Let's recall the relevant definitions and results from Olof's book:
\begin{itemize}
\item \textbf{Definition 9.5}. A linear system of curves of degree $d$ is a linear subvariety of the projective space of curves of degree $d$.

\item \textbf{Definition 9.10}. If
\[
D=m_1P_1+\cdots+m_rP_r
\]
is an effective zero-cycle on $P^2$, then $L_d(-D)$ denotes the linear system of plane curves of degree $d$ with multiplicity at least $m_i$ at $P_i$.

\item Also remember the convention
\[
\dim(\emptyset)=-1.
\]

\item \textbf{Definition 9.12}. The virtual dimension and expected dimension are given by
\[
\nu(L_d(-D))=\frac{d(d+3)}{2}-\sum_{i=1}^r \frac{m_i(m_i+1)}{2},
\qquad
e(L_d(-D))=\max\{-1,\nu(L_d(-D))\}.
\]

\item We say that
$
P_1,\dots,P_r
$
are in a general position if a certain linear map
\[
\phi:A_d\to k^s,
\]
obtained by evaluating polynomials and the relevant derivatives at the points, has maximal rank.

\item \textbf{Theorem 9.16}. If
$
P_1,\dots,P_r
$
are in general position and
$
D=P_1+\cdots+P_r,
$
then the dimension of $L_d(-D)$ is equal to the expected dimension.
\end{itemize}

\medskip

For
$
D_r=P_1+\cdots+P_r,$
by Definition 9.10 we have that
$
L_d(-D_r)
$
is the linear system of curves of degree $d$ passing through the points
$
P_1,\dots,P_r.
$
Since all multiplicities are equal to $1$, Definition 9.12 gives
\[
\nu(L_d(-D_r))=\frac{d(d+3)}{2}-r,
\qquad
e(L_d(-D_r))=\max\left\{-1,\frac{d(d+3)}{2}-r\right\}.
\]
Because the points are in general position, Theorem 9.16 applies exactly to
\[
D_r=P_1+\cdots+P_r,
\]
and therefore
\[
\dim(L_d(-D_r))=e(L_d(-D_r)).
\]

Therefore, whenever
\[
r=\frac{d(d+3)}{2},
\]
we obtain
\[
\dim(L_d(-D_r))=0.
\]
Since the convention is
\[
\dim(\emptyset)=-1,
\]
this shows that $L_d(-D_r)$ is nonempty.

By Definition 9.5, $L_d(-D_r)$ is a linear subvariety of the projective space of degree $d$ curves. A nonempty linear subvariety of projective dimension $0$ is $P^0$, and by Definition 8.1 with $n=0$, the space $P^0$ consists of exactly one point. Hence, if
\[
r=\frac{d(d+3)}{2},
\]
then there is exactly one degree $d$ curve passing through
$
P_1,\dots,P_r.
$

\medskip

\textbf{a)} For lines we have $d=1$. Then
\[
\frac{d(d+3)}{2}=\frac{1\cdot 4}{2}=2.
\]
Therefore
\[
\dim(L_1(-(P_1+P_2)))=0,
\]
so there is exactly one line passing through $P_1$ and $P_2$.

\medskip

\textbf{b)} For conics we have $d=2$. Then
\[
\frac{d(d+3)}{2}=\frac{2\cdot 5}{2}=5.
\]
Therefore
\[
\dim(L_2(-(P_1+\cdots+P_5)))=0,
\]
so there is exactly one conic passing through $P_1,\dots,P_5$.

\medskip

\textbf{c)} For cubics we have $d=3$. Then
\[
\frac{d(d+3)}{2}=\frac{3\cdot 6}{2}=9.
\]
Therefore
\[
\dim(L_3(-(P_1+\cdots+P_9)))=0,
\]
so there is exactly one cubic passing through $P_1,\dots,P_9$.

\medskip

\textbf{d)} For quartics we have $d=4$. Then
\[
\frac{d(d+3)}{2}=\frac{4\cdot 7}{2}=14.
\]
Thus $14$ points in general position are required in order to get a zero-dimensional linear system, and hence a unique quartic.

If $r<14$, then
\[
e(L_4(-D_r))=14-r>0,
\]
so Theorem 9.16 gives
\[
\dim(L_4(-D_r))>0.
\]
Therefore, there is a positive-dimensional family of quartics through those points, so the quartic is not uniquely determined.

Hence
\[
14
\]
points in a general position are required to specify a quartic.








\end{document}
